% Generated draft: requires mathematical and source review.
\documentclass[11pt]{article}
\usepackage[margin=25mm]{geometry}
\usepackage{fontspec}
\setmainfont{Latin Modern Roman}
\usepackage{amsmath,amssymb,mathtools,unicode-math}
\setmathfont{Latin Modern Math}
\usepackage{enumitem,needspace,xurl,hyperref}
\setlength{\parindent}{0pt}
\setlength{\parskip}{5pt}
\setlength{\emergencystretch}{3em}
\newcommand{\N}{\mathbb{N}}
\newcommand{\Z}{\mathbb{Z}}
\newcommand{\Q}{\mathbb{Q}}
\newcommand{\R}{\mathbb{R}}
\newcommand{\C}{\mathbb{C}}
\newcommand{\F}{\mathbb{F}}
\newcommand{\A}{\mathbb{A}}
\newcommand{\PP}{\mathbb{P}}
\newcommand{\E}{\mathbb{E}}
\newcommand{\eps}{\varepsilon}
\newcommand{\dd}{\,\mathrm{d}}
\begin{document}
\section*{UnsolvedMath ID: 30001077}
\textbf{Source problem code:} catalogue entry
\subsection{Definitions and mathematical statement}
Let \(\mu\) and \(\nu\) be Deligne--Mostow weight systems satisfying the source's integrality condition. Write \(U_\mu\) and \(U_\nu\) for the corresponding open moduli spaces of weighted configurations, and let \(F:U_\mu\to U_\nu\) be the stated forgetful map. In the customary setting these configurations are marked points on \(\PP^1\) modulo projective transformations, with weight-dependent complex-hyperbolic metrics; the precise weight and labeling conventions are not supplied.

Let \(\widehat U_\mu\) and \(\widehat U_\nu\) be the relevant metric completions. Determine when \(F\) extends to a map
\[
\overline F:\widehat U_\mu\longrightarrow\widehat U_\nu,
\qquad \overline F|_{U_\mu}=F.
\]
When it does, determine when the extension lifts to a map of the associated orbifolds. The source calls this equivalent to lifting to suitable torsion-free manifold covers. In that cover formulation the intended requirement is the existence of covers \(q_\mu:\widetilde U_\mu\to\widehat U_\mu\), \(q_\nu:\widetilde U_\nu\to\widehat U_\nu\), and a map \(\widetilde F\) such that
\[
q_\nu\circ\widetilde F=\overline F\circ q_\mu.
\]
The completion, regularity of the extended map, orbifold structures, and allowed covers must be specified before these criteria are precise. In particular, the asserted orbifold/cover equivalence is retained as source wording rather than proved for unspecified structures.
\subsection{Short English statement}
Classify when Deligne--Mostow forgetful maps extend across metric-completion strata and when those extensions respect orbifold structures or lift to suitable torsion-free covers. The weight, metric and covering conventions are missing from the extraction.
\subsection{Sources}
\begin{itemize}
\item Local UnsolvedMath registry, permanent ID \(30001077\), source code \texttt{catalogue entry}, ``Orbifold Lifts of Deligne Mostow Forgetful Maps.'' Registry curation: Ulam AI / UnsolvedMath Contributors, CC BY 4.0; underlying sources retain their supplied provenance.
\item \emph{Geometry and Arithmetic around Hypergeometric Functions} (2009), Oberwolfach Reports, pp.\ 2557--2620, DOI \url{https://doi.org/10.4171/owr/2008/45}; supplied label ``OWR Report 45.''
\item Supplied literature triage, \texttt{literature\_checked\_at}: 2026-08-21: labeled open; says the passage seeks a general extension criterion for the relevant hypergeometric orbifold maps and no resolving result for the general existence and lifting question was found. This is supplied dated curation, not independently verified.
\end{itemize}
\end{document}
